% ===========================================================================
% Motivating Example -- paper-ready LaTeX
% Pairs with Figure 1 (motivation_figure.tex, Skill `d3-visualization`).
% The figure puts the document on the left and the code on the right and marks
% each inconsistency with a number at both of its ends.
% Replace \ref{sec:corpus}, \ref{sec:taxonomy}, \ref{fig:motivation}.
% ===========================================================================
\section{Motivating Example}
\label{sec:motivation}
%
% Local only: long unbreakable \texttt{} tokens would otherwise overflow.
% Verified clean from 8cm to 16cm text width.
\begingroup\emergencystretch=3em

An Agent Skill is a directory an agent loads as instructions. It has two sides:
a \texttt{SKILL.md} contract written for the reader, and the scripts that are
supposed to carry it out. The agent cannot execute the contract to see whether
it holds. It reads the text and acts on it. We take one Skill from the SkillBench
corpus we audit in Section~\ref{sec:corpus} and read its contract the way an
agent would.

Figure~\ref{fig:motivation} puts the document on the left and the code on the
right, and numbers every inconsistency at both of its ends. The boxes those ends
sit in name its boundary. Both ends on the left is a document--document
inconsistency; one end on each side is a document--code inconsistency.

The first defect has both ends in the document. Line~3 of \texttt{SKILL.md}
describes the Skill as producing visualizations ``with D3.js (v6)''. Line~75
instructs the reader to pin D3 to ``a specific version (default:
\texttt{d3@7.9.0})'', and the next line recommends vendoring
\texttt{d3.v7.9.0.min.js}. One file, two major versions of its own central
dependency.

The second and third cross from the document into the code, and both are
determinism rules that the shipped example breaks. Lines~59--60 require the
reader to ``sort input rows deterministically before binding to marks'' and to
use a ``stable grouping order (explicit
\texttt{Array.from(grouped\allowbreak{}.keys()).sort()})''.
\texttt{bubble\_\allowbreak{}chart\_\allowbreak{}example.js:62} builds its sector grouping as
\texttt{[...new\allowbreak{}Set(data.map(d\allowbreak{} => d.sector))]}, which keeps the input order. The
Skill's other example does sort its rows
(\texttt{interactive\_\allowbreak{}table\_\allowbreak{}example.js:73}), so the two shipped examples
disagree with each other as well.

The third defect is the same shape. Lines~68--71 state that if a force layout is
required the reader must ``fix the tick count, fix initial positions
deterministically'' and ``run exactly N ticks and stop''.
\texttt{bubble\_\allowbreak{}chart\_\allowbreak{}example.js:56} builds a default
\texttt{d3.force\allowbreak{}Simulation(data)}: a tick count nobody sets, initial positions
nobody fixes, and no seed. An agent that copies the example inherits a layout the
contract forbids.

All three are settled by opening the two boxes, before any agent runs, and all
three are instances of the types we define in Section~\ref{sec:taxonomy}.

\medskip\noindent\textbf{Observation 1.} \emph{The boundary is a property of the
artifacts, not a label we assign.} The version is contradicted inside one file;
the two determinism rules are contradicted across two. Neither is a
documentation-style problem: link checkers, linters and doc builds all pass on
this Skill. RQ1 asks which kinds of contract inconsistency occur in real Skills,
how often, and how they distribute across the boundaries.

\medskip\noindent\textbf{Observation 2.} \emph{Reading the artifacts tells you
what is inconsistent, not what an agent will do about it.} Both determinism
rules are stated as requirements, and both are broken silently: an unsorted
grouping and a default force layout both draw a perfectly plausible chart, only
not the same one twice. An agent may follow the rules, ignore them, or never read
them, and the artifacts do not say which. RQ2 asks whether these inconsistencies
can be found automatically, and whether the ones an agent will build on can be
separated from the ones it will read past.

\medskip\noindent\textbf{Observation 3.} \emph{They are not equally repairable.} Correcting the version is a one-token edit once someone decides
which line to trust. Correcting either determinism rule means changing a layout
or a sort and accepting that the picture will move, and the two fixes land in
different places: one in the grouping code, one in the simulation setup. Repair
therefore has to be evaluated per defect, not per Skill. RQ3 asks what each kind
of inconsistency costs an agent, and whether repairing it recovers that cost.
\endgroup
